\section{Metric--Measure 확장}

기하 층에는 불확실하거나 중복된 entity가 포함될 수 있으므로 set-level
correspondence만으로는 충분하지 않다. 유한 metric carrier에 확률질량을
부여하고 coupling을 통해 질량을 비교한다. 아래 결과는 exactness theorem이며
finite-sample convergence theorem이 아니다. Coupling과 Gromov--Wasserstein의
배경은 표준 이론을 따른다
\cite{memoli2011gromov,villani2009optimal,peyre2019computational}.

\begin{definition}[완전지지 metric--measure 상태]
완전지지 metric--measure 상태는 $M=(X,d,\mu,\lambda)$이다. 여기서
$(X,d)$는 유한 metric space, $\mu:X\to(0,1]$는
$\sum_{x\in X}\mu(x)=1$을 만족하는 확률질량, $\lambda:X\to\mathcal L$는
label map이다. 완전지지는 모든 $x\in X$에 대해 $\mu(x)>0$임을 뜻한다.
\end{definition}

완전지지는 exactness 결과에 필요하다. 질량이 0인 entity는 coupling objective를
바꾸지 않고 추가되거나 이동될 수 있기 때문이다. 측도는 상태의 일부이지만
좌표나 목록에 없는 attribute를 기하학적 증거로 암묵적으로 취급하지 않는다.

\begin{definition}[label-compatible coupling]
$M_X=(X,d_X,\mu_X,\lambda_X)$와
$M_Y=(Y,d_Y,\mu_Y,\lambda_Y)$에 대해 coupling은
$\pi:X\times Y\to\R_{\geq0}$ 중
\[
 \sum_{y\in Y}\pi(x,y)=\mu_X(x),
 \qquad
 \sum_{x\in X}\pi(x,y)=\mu_Y(y)
\]
를 만족하는 함수이다. $\pi(x,y)>0$이면
$\lambda_X(x)=\lambda_Y(y)$일 때 label-compatible이라 한다.
\end{definition}

\begin{definition}[$p$-metric--measure discrepancy]
$1\leq p<\infty$와 label-compatible coupling $\pi$에 대해
\[
 \mathcal D_p(\pi)^p=
 \sum_{x,x'\in X}\sum_{y,y'\in Y}
 |d_X(x,x')-d_Y(y,y')|^p
 \pi(x,y)\pi(x',y')
\]
로 정의한다. 그리고
\[
 \Delta_{\mathrm{mm},p}(M_X,M_Y)=\inf_{\pi}\mathcal D_p(\pi)
\]
로 둔다. infimum은 label-compatible coupling 전체에 대해 취하며, 그러한
coupling이 없으면 discrepancy를 $+\infty$로 둔다.
\end{definition}

이는 hard label constraint를 갖는 유한 Gromov--Wasserstein-type objective이다
\cite{memoli2011gromov}.
좌표 비용, soft attribute 비용, 학습된 correspondence는 포함하지 않는다.
예를 들어 feature--structure joint cost는 fused Gromov--Wasserstein 정식화로
이어진다 \cite{vayer2020fused}. 그러한 항을 추가하면 별도의 discrepancy가 되며
식별가능성을 별도로 분석해야 한다.

\begin{theorem}[Metric--measure discrepancy 0의 exactness]
$M_X,M_Y$가 완전지지 metric--measure 상태라고 하자. 모든
$1\leq p<\infty$에 대해
$\Delta_{\mathrm{mm},p}(M_X,M_Y)=0$일 필요충분조건은
bijection $\varphi:X\to Y$가 존재하여 모든 $x,x'\in X$에 대해
\[
 d_Y(\varphi(x),\varphi(x'))=d_X(x,x'),\qquad
 \lambda_Y(\varphi(x))=\lambda_X(x),\qquad
 \mu_Y(\varphi(x))=\mu_X(x)
\]
를 만족하는 것이다.
\end{theorem}

\begin{proof}
먼저 그러한 bijection이 존재한다고 하자. 다음과 같이 둔다.
\[
 \pi_\varphi(x,y)=
 \begin{cases}
 \mu_X(x),&y=\varphi(x),\\
 0,&\text{otherwise}.
 \end{cases}
\]
첫 번째 marginal은 구성상 $\mu_X$이다. $y=\varphi(x)$일 때 bijection과
measure preservation에 의해 두 번째 marginal은
$\mu_X(x)=\mu_Y(y)$가 된다. 따라서 이는 coupling이고, label 보존에
의해 label-compatible이다. 양의 product가 나타나는 항에서는
$y=\varphi(x)$와 $y'=\varphi(x')$이며 거리 차이는 0이다. 따라서
$\mathcal D_p(\pi_\varphi)=0$이다.

반대로 discrepancy가 0이라고 하자. 모든 coupling의 집합은 유한차원 공간
$\R^{X\times Y}$에서 closed and bounded이며, 유한한 label constraint를
적용한 부분집합도 compact이다. coupling이 존재하면 $\mathcal D_p$의
연속성으로 최소 coupling $\pi_*$가 존재하고
$\mathcal D_p(\pi_*)=0$이다.

$C=\{(x,y):\pi_*(x,y)>0\}$라 하자. 양쪽 marginal의 full support로부터
모든 $x$와 모든 $y$가 $C$의 어떤 pair에 포함되므로 $C$는 correspondence이다.
$C$의 두 pair $(x,y),(x',y')$에 대해 정의식의 계수는 양수이다. 모든
summand가 음이 아니고 합이 0이므로
$d_X(x,x')=d_Y(y,y')$이다.

$(x,y),(x,y')$가 모두 $C$에 있으면 위 등식으로
$d_Y(y,y')=d_X(x,x)=0$이므로 $y=y'$이다. 같은 방식으로
$(x,y),(x',y)$가 모두 $C$에 있으면 $x=x'$이다. 따라서 $C$는 bijection
$\varphi:X\to Y$의 graph이다. 거리 등식은 isometry를, positive support는
label preservation을 준다. 마지막으로 첫 번째 marginal은
$\mu_X(x)=\pi_*(x,\varphi(x))$를 주고, graph 성질과 두 번째 marginal은
$\mu_Y(\varphi(x))=\pi_*(x,\varphi(x))$를 준다. 따라서 bijection은
measure-preserving이다.
\end{proof}


이 정리는 metric--measure 층의 zero set을 정확히 해석한다. 그러나 sampling
robustness, empirical convergence rate, 임의 예측 담체에서 좋은 coupling의
존재, coordinate loss에 대한 우월성은 증명하지 않는다. 이는 별도의 경험적
주장이다.


\begin{proposition}[Label-compatible coupling의 feasibility]
Label-compatible coupling이 존재하는 필요충분조건은 모든 label의 total mass가
일치하는 것이다.
\[
\sum_{x:\lambda_X(x)=\ell}\mu_X(x)
=
\sum_{y:\lambda_Y(y)=\ell}\mu_Y(y)
\qquad(\ell\in\mathcal L).
\]
\end{proposition}

\begin{proof}
필요성은 앞의 marginal 계산에서 따른다. 충분성을 위해 각 label $\ell$의
두 label block 안에서 common mass로 정규화한 product coupling을 취하고,
이를 모든 label에 대해 합한다. 그러면 요구된 marginal을 가지며 cross-label
mass가 없는 coupling을 얻는다. Common mass가 0인 block은 생략한다.
\end{proof}

Tagged label map을 $x\in F(A)$에 대해 $\bar\lambda(x)=(A,\lambda_A(x))$로,
$y\in G(A)$에 대해 $\bar\nu(y)=(A,\nu_A(y))$로 정의한다.

\begin{corollary}[State-level metric--measure invariance]
$\Sigma\cong_{\mathrm{str}}\Sigma'$이면 induced full-support metric--measure
state $(|F|,d,\mu,\bar\lambda)$와 $(|G|,\rho,\eta,\bar\nu)$의 metric--measure
discrepancy는 모든 $1\leq p<\infty$에 대해 0이다.
\end{corollary}

\begin{proof}
Structural isomorphism을 measure-preserving, label-preserving isometry로
사용하고 zero metric--measure discrepancy theorem을 적용한다.
\end{proof}
\begin{proposition}[Hard-label discrepancy의 extended pseudometric 성질]
Tagged label을 갖는 full-support finite metric--measure state에서
$\Delta_{\mathrm{mm},1}$은 symmetric이고 diagonal에서 0이며 triangle
inequality를 만족한다. Label-compatible coupling이 없으면 값은 $+\infty$일
수 있다.
\end{proposition}

\begin{proof}
Identity-graph coupling이 diagonal vanishing을 준다. Label-compatible
coupling을 transpose하면 marginal 순서가 뒤바뀌고 $\mathcal D_1$ 값은
그대로이므로 infimum을 취해 symmetry를 얻는다. Triangle inequality를 위해
label-compatible coupling $\pi_{XY}$와 $\pi_{YZ}$를 택한다. Full support로
$\mu_Y(y)>0$이므로
\[
\gamma(x,y,z)=\frac{\pi_{XY}(x,y)\pi_{YZ}(y,z)}{\mu_Y(y)}
\]
는 well-defined이다. 이는 optimal transport의 표준 gluing construction이다
\cite{villani2009optimal,peyre2019computational}. 직접 합하면 $(X,Y)$와 $(Y,Z)$ marginal이 각각
원래 coupling임을 알 수 있다. $\pi_{XZ}$를 그 $(X,Z)$ marginal이라 하자.
$\gamma(x,y,z)>0$이면
$\lambda_X(x)=\lambda_Y(y)=\lambda_Z(z)$이므로 $\pi_{XZ}$도
label-compatible이다. $\gamma$에서 뽑은 두 triple에 대해
\[
|d_X(x,x')-d_Z(z,z')|
\leq |d_X(x,x')-d_Y(y,y')|+|d_Y(y,y')-d_Z(z,z')|
\]
이다. 이를 $\gamma\otimes\gamma$에 대해 합하면
$\mathcal D_1(\pi_{XZ})\leq\mathcal D_1(\pi_{XY})+
\mathcal D_1(\pi_{YZ})$를 얻는다. 두 원래 coupling에 대해 infimum을 취하면 triangle inequality가
따른다. 어느 한 coupling이 infeasible이면 대응하는 extended value가
$+\infty$이므로 부등식은 자명하다.
\end{proof}
